% !TeX program = xelatex
\documentclass[11pt]{impeart}

% Das Handbuch verwendet bewusst einen globalen Hausstil. Gewöhnliche Dokumente
% sollten normalerweise \UseFont ohne Modus oder den Template-Schlüssel `fonts` verwenden.
\UseFont{libertinus}[global]
\setmonofont[
  BoldFont=lmmonolt10-bold.otf,
  ItalicFont=lmmono10-italic.otf,
  BoldItalicFont=lmmonolt10-boldoblique.otf
]{lmmono10-regular.otf}
\UseFeatures{hyperlinks,tables}

\usepackage{catchfile}
\usepackage{metalogo}
\usepackage{pdfpages}
\usepackage[ngerman]{babel}

\newcommand{\IMPEVersionFile}{../../VERSION}
\newcommand{\IMPEShowcaseFile}{../../_showcase/main.pdf}

\IfFileExists{\IMPEVersionFile}
  {\CatchFileEdef{\IMPEVersion}{\IMPEVersionFile}{\endlinechar=-1\relax}}
  {\PackageError{impe-manual}
     {Erforderliche Repository-Ressource ../../VERSION fehlt}
     {Aus doc/de kompilieren oder scripts/build_manual.ps1 aus dem Repository-Stammverzeichnis ausfuehren.}}

% Das erzeugte Handbuch bleibt reproduzierbar, wenn SOURCE_DATE_EPOCH festgelegt ist.
\ifdefined\XeTeXversion
  \special{pdf:trailerid [
    <5677bc84f2edee3de328d781062f390b>
    <778a1a55595b871c7ecbd03469f75deb>
  ]}
\else\ifdefined\pdftrailerid
  \expandafter\pdftrailerid\expandafter{IMPE-\IMPEVersion-manual-de}
\fi
\fi

\title{IMPE \LaTeX{} System}
\subtitle{Benutzerhandbuch}
\author{YANG Sikai}
\date{Version \IMPEVersion}

\newcommand{\IMPERepo}{https://github.com/KelvinYangBK67/IMPE-LaTeX-System}
\newcommand{\IMPEIssues}{https://github.com/KelvinYangBK67/IMPE-LaTeX-System/issues}
\newcommand{\IMPEShowcaseURL}{https://github.com/KelvinYangBK67/IMPE-LaTeX-System/blob/main/_showcase/main.pdf}
\newcommand{\IMPETaggedShowcaseURL}{https://github.com/KelvinYangBK67/IMPE-LaTeX-System/blob/v\IMPEVersion/_showcase/main.pdf}

\begin{document}

\maketitle

\begin{center}
\small Deutsche Übersetzung der englischen Referenzausgabe
\end{center}

\tableofcontents

\section{Einführung}

IMPE steht für \emph{Integrated Multilingual Publishing Environment}.
Das IMPE \LaTeX{} System ist ein modulares Dokumentsystem für \XeLaTeX{}, das
für Arbeitsabläufe entwickelt wurde, in denen Dokumentlayout, optionale
Funktionalität, Schriftauswahl und mehrsprachiger Satz wiederverwendbar sein
sollen, anstatt im Vorspann jedes Dokuments neu aufgebaut zu werden. Der
ausgeschriebene Name beschreibt das Hauptziel des Systems: eine integrierte
Publikationsumgebung für Dokumente bereitzustellen, die mehrere Sprachen,
Schriftsysteme, Layouts und Funktionsmodule miteinander verbinden können.

Das System entstand aus mehrsprachiger Dokumentarbeit, bei der ein
konventioneller Vorspann schnell schwer wartbar wurde. Ein Dokument kann ein
stabiles Seitendesign, eine lateinische Textschrift, CJK-Unterstützung, ein oder
mehrere historische Schriftsysteme, schriftspezifisches Shaping,
Bibliographieunterstützung, Bilder, Tabellen sowie unterschiedliche
Einstellungen für Artikel, Bücher und Folien benötigen. Keine dieser Aufgaben
ist für sich genommen ungewöhnlich, doch ihre Kombination führt leicht zu
wiederholtem Einrichtungscode und einer wachsenden Zahl dokumentspezifischer
Ausnahmen.

IMPE begegnet diesem Problem, indem das System in wiederverwendbare Teile
aufgeteilt wird. Aus Sicht des Benutzers ist die wichtigste Unterscheidung die
zwischen \emph{Layouts}, \emph{Schriften} und \emph{Features}. Ein Layout
beschreibt die grundlegende visuelle Organisation eines Dokuments.
Schriftfamilien stellen globales oder lokales Schriftverhalten bereit,
einschließlich schriftsystembewussten Routings, wo dies erforderlich ist.
Features stellen optionale Funktionen wie Mathematik, Hyperlinks, Zitate,
Tabellen oder Bilder bereit. Diese Teile können unabhängig voneinander
ausgewählt oder über eine einheitliche Template-Schnittstelle kombiniert
werden.

Die normale Verwendung von IMPE ist bewusst knapp gehalten. Eine Wrapper-Klasse
kann ein sinnvolles Dokumentlayout und eine globale Standardschriftkonfiguration
festlegen; anschließend lädt der Benutzer nur die zusätzlichen Familien oder
Features, die das Dokument benötigt. Wenn eine explizitere Kontrolle gewünscht
ist, lassen sich dieselben Komponenten in einem einzigen Aufruf von
\verb|\UseTemplateSet| deklarieren.

\XeLaTeX{} ist die primäre Dokument-Engine für IMPE. Diese Wahl ermöglicht es
dem System, sich auf Unicode-Text, OpenType-Schriften, \texttt{fontspec} und die
Inter-Character-Mechanismen von \XeTeX{} zu stützen und zugleich das vertraute
\LaTeX{}-Dokumentmodell beizubehalten. Einige spezialisierte interne Routen
können andere Engines aufrufen, wenn ein bestimmtes Schriftsystem oder ein
bestimmter Renderer dies erfordert; gewöhnliche IMPE-Dokumente sind jedoch zur
Kompilierung mit \XeLaTeX{} vorgesehen.

Dieses Handbuch ist eine benutzerorientierte Einführung und keine erschöpfende
API-Referenz. Detaillierte Dokumentation der Subsysteme wird separat im
Repository unter \texttt{docs/} gepflegt. Insbesondere die Dokumente zu
Schriften, Layouts, Features und Systemarchitektur beschreiben
Registrierungsfelder und Implementierungsdetails, deren Wiederholung hier
unnötig wäre. Das englische Handbuch ist die Referenzausgabe; übersetzte
Handbücher sollen seiner Struktur und seinem technischen Inhalt folgen.

Eine visuelle Übersicht über das System bietet der Projekt-Showcase unter
\href{\IMPEShowcaseURL}{dem aktuellen IMPE-Showcase-PDF}. Die mit diesem
Handbuch eingefrorene Ausgabe ist über den
\href{\IMPETaggedShowcaseURL}{Versions-Tag} verfügbar. Dieses kanonische
Repository-PDF wird außerdem in Anhang~\ref{app:showcase} wiedergegeben. Es soll
die Bandbreite an Schriftsystemen, Schreibrichtungen, Layouts und
Feature-Kombinationen demonstrieren, aus denen das System hervorgegangen ist.

\section{Installation}

\subsection{Distributionsformen}

IMPE wird in drei Distributionsformen vorbereitet, die jeweils für einen
anderen Anwendungsfall vorgesehen sind.

Die \emph{CTAN-orientierte Distribution} enthält die portable Laufzeitumgebung
und die öffentliche Dokumentation, jedoch nicht die private oder lokal
gepflegte Schriftbibliothek. Diese Form ist für die reguläre
\TeX{}-Distribution vorgesehen. Die \emph{Core-Distribution} enthält dieselbe
Systemlogik für Benutzer, die IMPE lieber direkt aus einem Release installieren
und ihre Schriften separat verwalten. Die \emph{Full-Distribution} kann eine
vorbereitete lokale Schriftbibliothek hinzufügen, soweit der
Weiterverteilungsstatus der einzelnen Schriften dies erlaubt.

Diese Unterscheidung ist beabsichtigt. IMPE selbst ist ein Dokumentsystem; eine
große Schriftsammlung gehört nicht zum logischen Kern des Pakets. Die Trennung
macht die Laufzeitumgebung portabel und erlaubt verschiedenen Rechnern oder
Benutzern, unterschiedliche Schriftbibliotheken bereitzustellen, ohne das
Framework zu verändern.

\subsection{Installationsort}

Eine direkte Installation legt das Paket im TEXMF-Baum des Benutzers unter

\begin{verbatim}
tex/latex/impe/
\end{verbatim}

ab, sodass Paket und Wrapper-Klassen über den normalen \TeX{}-Suchpfad
verfügbar sind. Die kanonischen öffentlichen Einstiegspunkte sind

\begin{verbatim}
impe.sty
impeart.cls
impebook.cls
impereport.cls
impebeamer.cls
\end{verbatim}

mit entsprechenden \texttt{\_zh}-Wrapper-Klassen für chinesisch orientierte
Dokumentstandards. Kompatibilitätseinstiege unter den früheren
\texttt{next*}-Namen werden ebenfalls installiert; siehe
Abschnitt~\ref{sec:compatibility}.

Bei Verwendung eines direkten IMPE-Releases kann der mitgelieferte Installer
die Laufzeitumgebung in den TEXMF-Baum des Benutzers kopieren. Ein
Release-Archiv kann den üblichen Windows-Einstiegspunkt

\begin{verbatim}
install.bat
\end{verbatim}

bereitstellen. Der Launcher delegiert an \texttt{install.ps1}; fortgeschrittene
Benutzer können dieses Skript mit dem üblichen Skript-Runner ihres Systems
direkt aufrufen.

Benutzer, die über eine \TeX{}-Distribution installieren, benötigen den
Repository-seitigen Installer nicht; die Paketerkennung übernimmt die
\TeX{}-Distribution auf dem üblichen Weg.

\subsection{Schriftressourcen und lokale Konfiguration}

Das Quell-Repository verfolgt nicht die vollständige lokale Schriftbibliothek.
Bei einer lokalen oder Full-Installation werden katalogisierte Schriftdateien
normalerweise unterhalb eines Stammverzeichnisses \texttt{assets/fonts/}
aufgelöst. Ein anderer Ort kann gewählt werden, ohne das Paket selbst zu
bearbeiten.

Der bevorzugte dauerhafte Mechanismus ist \texttt{impe.local.tex}. Eine lokale
Konfiguration kann beispielsweise Folgendes enthalten:

\begin{verbatim}
\SetCatalogFontRoot{D:/Fonts/IMPE}
\end{verbatim}

wobei der Pfad an den lokalen Rechner angepasst wird. Derselbe Befehl kann auch
aus einem Dokument heraus verwendet werden, wenn eine einmalige Überschreibung
geeigneter ist. Core- und CTAN-Installationen bleiben ohne private
Schriftsammlung verwendbar; nur Dokumentfunktionen, die nicht verfügbare lokale
Schriftdateien anfordern, benötigen naturgemäß die entsprechenden Ressourcen.

\subsection{Schneller Installationstest}

Nach der Installation sollte sich die folgende Datei auf einem System mit
verfügbarem Libertinus mit \XeLaTeX{} kompilieren lassen:

\begin{verbatim}
\documentclass{impeart}
\UseFont{libertinus}

\begin{document}
\LIB{IMPE is ready.}
\end{document}
\end{verbatim}

Wenn dies gelingt, sind die kanonische Wrapper-Klasse, der Paketeinstieg und die
automatische Schnittstelle zum Laden von Schriften für \TeX{} sichtbar.

\section{Schnellstart}

Der einfachste IMPE-Arbeitsablauf besteht darin, eine Wrapper-Klasse zu wählen,
die gewünschten Schriftfamilien oder Features auszuwählen und das Dokument
anschließend wie gewohnt zu schreiben. IMPE ersetzt die gewöhnliche
\LaTeX{}-Dokumentstruktur nicht, sondern legt eine wiederverwendbare
Einrichtungsschicht darum.

Ein kleiner Artikel kann daher beispielsweise so aussehen:

\begin{verbatim}
\documentclass[11pt]{impeart}

\UseFont{libertinus}
\UseFeatures{math,hyperlinks}

\title{A Sample Document}
\subtitle{Typeset with IMPE}
\author{Author Name}

\begin{document}

\maketitle

\section{Introduction}
\LIB{Hello, world.}

\end{document}
\end{verbatim}

Die Klasse liefert bereits das normale englische Artikellayout und die
Basisschrift. Der Aufruf \verb|\UseFont{libertinus}| ohne Modus folgt dem
registrierten automatischen Verhalten dieser Familie, aktiviert ihren globalen
Schriftstapel und stellt zugleich den lokalen Befehl \verb|\LIB| bereit,
während \verb|\UseFeatures| zwei optionale Feature-Module aktiviert. Der Rest
ist gewöhnliches \LaTeX{}.

Für ein Dokument, dessen Einrichtung sich besser als einzelne Deklaration
ausdrücken lässt, kann dieselbe Idee über \verb|\UseTemplateSet| geschrieben
werden:

\begin{verbatim}
\documentclass{impeart}

\UseTemplateSet{
  fonts = {libertinus},
  features = {math,hyperlinks}
}
\end{verbatim}

Ein Template-Set kann außerdem ein Layout und zusätzliche Schriftfamilien
auswählen. In der Praxis liefern Wrapper-Klassen bereits ein geeignetes
Standardlayout, sodass die meisten Dokumente nur jene Einstellungen angeben,
die davon abweichen.

IMPE kann auch als Paket unterhalb einer standardmäßigen \LaTeX{}-Klasse
geladen werden:

\begin{verbatim}
\documentclass{article}
\usepackage{impe}

\UseTemplateSet{
  layout = en_doc,
  fonts = {cmu,libertinus},
  features = {hyperlinks}
}
\end{verbatim}

Diese Form ist nützlich, wenn die zugrunde liegende Klasse ausdrücklich
beibehalten werden muss oder wenn jede IMPE-Komponente manuell ausgewählt werden
soll. Wrapper-Klassen und das rohe Paket verwenden dasselbe zugrunde liegende
System.

\section{Einsatzbereite Vorlagen}

Das Repository enthält unter \texttt{templates/} eine kleine Sammlung von
Starterprojekten. Sie sind für das tatsächliche Schreiben vorgesehen und nicht
für Regressionstests oder interne Demonstrationen. Die Dateien unter
\texttt{examples/} dienen dagegen in erster Linie dem Debugging und der
Überprüfung.

Die derzeitigen Startervorlagen sind:

\begin{description}
  \item[\texttt{article\_en/}] Englische Artikelvorlage mit \texttt{impeart}.
  \item[\texttt{article\_zh/}] Chinesische Artikelvorlage mit \texttt{impeart\_zh}.
  \item[\texttt{book\_en/}] Englische Buchvorlage mit \texttt{impebook}.
  \item[\texttt{book\_zh/}] Chinesische Buchvorlage mit \texttt{impebook\_zh}.
  \item[\texttt{beamer\_en/}] Englische Präsentationsvorlage mit \texttt{impebeamer}.
  \item[\texttt{beamer\_zh/}] Chinesische Präsentationsvorlage mit \texttt{impebeamer\_zh}.
\end{description}

Jede Vorlage folgt dem Stil eines installierten Pakets und enthält realistische
Metadaten, Gliederung und Dokumentinhalt sowie gegebenenfalls repräsentatives
Tabellen-, Abbildungs- oder Folienmaterial. Ein neues Projekt kann daher damit
beginnen, das nächstgelegene Vorlagenverzeichnis zu kopieren und den
Dokumentinhalt zu bearbeiten, anstatt einen Vorspann von Grund auf neu
aufzubauen.

Die Vorlagen verwenden bewusst die kanonischen \texttt{impe*}-Einstiegspunkte.
Sie zeigen die normale benutzerseitige Praxis, während Repository-lokale
\texttt{examples/} Pfade oder Fixtures verwenden können, die ausschließlich
für Entwicklung und Regressionstests bestimmt sind.

\section{Dokumentklassen und öffentliche Einstiegspunkte}

\subsection{Wrapper-Klassen}

IMPE stellt Wrapper-Klassen für die vier wichtigsten Dokumentformen des Systems
bereit:

\begin{description}
  \item[\texttt{impeart}] Englisch orientierter Artikel-Wrapper. Das normale
    Layout ist \texttt{en\_doc}, die globale Standardfamilie ist \texttt{cmu}.
  \item[\texttt{impebook}] Englisch orientierter Buch-Wrapper. Das normale
    Layout ist \texttt{en\_book}, die globale Standardfamilie ist \texttt{cmu}.
  \item[\texttt{impereport}] Englisch orientierter Report-Wrapper. Er verwendet
    standardmäßig das Layout \texttt{en\_doc} und \texttt{cmu}.
  \item[\texttt{impebeamer}] Präsentations-Wrapper. Er wählt standardmäßig das
    Layout \texttt{beamer} und \texttt{cmu}.
\end{description}

Chinesisch orientierte Gegenstücke werden als \texttt{impeart\_zh},
\texttt{impebook\_zh}, \texttt{impereport\_zh} und
\texttt{impebeamer\_zh} bereitgestellt. Sie wählen die entsprechenden
chinesischen Dokumenteinstellungen und nehmen neben \texttt{cmu} auch
\texttt{shanggu} in die globale Standardschriftkonfiguration auf. Die
Unterscheidung betrifft die Standardkonfiguration und schränkt nicht ein, welche
Sprachen im Dokument vorkommen dürfen.

Unbekannte Klassenoptionen werden an die zugrunde liegende Standardklasse
weitergereicht. Eine normale Option wie

\begin{verbatim}
\documentclass[12pt,twoside]{impeart}
\end{verbatim}

steuert daher weiterhin wie gewohnt die zugrunde liegende Klasse
\texttt{article}.

\subsection{\texttt{impe.sty} direkt verwenden}

Die Wrapper-Klassen sind Komfortschnittstellen und keine separaten
Implementierungen. Ein Benutzer kann stattdessen eine Standardklasse
beibehalten und

\begin{verbatim}
\usepackage{impe}
\end{verbatim}

laden, gefolgt von einer expliziten Konfiguration für Layout, Schriften und
Features. Dies ist der empfohlene Weg, wenn die Dokumentklasse von einer
Zeitschrift, einem Verlag oder einem anderen Paket vorgegeben wird und nicht
durch einen IMPE-Wrapper ersetzt werden soll.

\subsection{Titel und Untertitel}

IMPE unterstützt \verb|\subtitle{...}| neben den Standardbefehlen
\verb|\title| und \verb|\author|. Wrapper-Klassen wenden ihre eigenen
Standardwerte für den Titelblock an, während der gewöhnliche
\LaTeX{}-Titelmechanismus erkennbar bleibt. Ein typischer Vorspann lautet
einfach:

\begin{verbatim}
\title{Main Title}
\subtitle{A More Specific Description}
\author{Author Name}
\end{verbatim}

gefolgt von \verb|\maketitle| im Dokumentkörper.

\section{Das Template-Modell}

\subsection{Vier Konfigurationsschlüssel}

Der wichtigste einheitliche Befehl ist \verb|\UseTemplateSet|. Sein
öffentliches Konfigurationsmodell ist bewusst klein gehalten:

\begin{verbatim}
\UseTemplateSet{
  layout = <preset>,
  fonts = {<families>},
  globalfonts = {<families>},
  features = {<features>}
}
\end{verbatim}

Die vier Schlüssel haben unterschiedliche Aufgaben.

\begin{description}
  \item[\texttt{layout}] wählt eine öffentliche Layout-Voreinstellung, die zur
    aktiven Dokumentklasse passt.
  \item[\texttt{fonts}] lädt Familien über ihr registriertes automatisches
    Verhalten. Dies ist die normale Wahl: Eine Familie kann einen lokalen
    Befehl definieren, eine dokumentweite Schrift anwenden oder
    bereichsspezifisches Routing einrichten.
  \item[\texttt{globalfonts}] zwingt die aufgeführten Familien ausdrücklich
    durch ihren globalen Modus oder ihren globalen Routing-Pfad.
  \item[\texttt{features}] lädt optionale Funktionsmodule.
\end{description}

\texttt{mainfonts} wird als Alias von \texttt{globalfonts} akzeptiert; beide
Schlüssel sind explizite globale Überschreibungen und nicht die normale
automatische Schnittstelle.

Die Template-Schnittstelle verlangt nicht, dass jedes Dokument alle vier
Schlüssel deklariert. Ihr Zweck besteht darin, einen vorhersehbaren Ort
bereitzustellen, an dem eine vollständige Konfiguration \emph{ausgedrückt
werden kann}. Wrapper-Klassen wenden bereits ihre eigenen Layout- und globalen
Schriftstandards an; nur eine zusätzliche Feature-Liste anzugeben ist daher
vollkommen normal.

\subsection{Spezialisierte Kurzbefehle}

Für Dokumente, die keine vollständige Template-Deklaration benötigen, bietet
IMPE öffentliche Kurzbefehle für die einzelnen Subsysteme:

\begin{verbatim}
\UseLayout{en_doc}
\UseLayouts{...}

\UseFont{libertinus}
\UseFonts{arabic,tibetan}

\UseFont{libertinus}[global]
\UseLocalFonts{arabic,tibetan} % explicit local override
\UseGlobalFonts{libertinus}    % explicit global override

\UseFeature{hyperlinks}
\UseFeatures{math,hyperlinks,tables}
\end{verbatim}

Die Font-Formen ohne Modus sind die empfohlenen Kurzbefehle für den normalen
Gebrauch. Die lokalen und globalen Aliase bleiben nützlich, wenn ein Dokument
das registrierte Verhalten einer Familie bewusst überschreibt. Die
vereinheitlichte und die spezialisierte Schnittstelle sind zwei Ansichten
desselben öffentlichen Modells und keine konkurrierenden
Konfigurationssysteme.

\subsection{Komposition statt monolithischer Templates}

Ein IMPE-Template ist bewusst keine unteilbare Stildatei. Ein Layout kann mit
unterschiedlichen Schriftauswahlen wiederverwendet werden, dieselbe
Feature-Menge kann in einem Artikel oder Buch zum Einsatz kommen, und eine
schriftspezifische Schriftfamilie kann ergänzt werden, ohne nicht verwandte
Dokumentfunktionen zu verändern. Diese Komposition ist das zentrale
Gestaltungsprinzip der öffentlichen Schnittstelle.

Die folgenden Deklarationen beschreiben beispielsweise zwei sehr
unterschiedliche Dokumente, ohne zwei voneinander unabhängige
Vorspann-Templates zu erfordern:

\begin{verbatim}
% English article
\UseTemplateSet{
  layout = en_doc,
  fonts = {libertinus},
  features = {math,hyperlinks}
}

% Chinese book with additional script support
\UseTemplateSet{
  layout = zh_book,
  fonts = {cmu,shanggu,sanskrit,tibetan},
  features = {citations,index,hyperlinks}
}
\end{verbatim}

Das zweite Beispiel nennt nur die Familien und Features, die es benötigt; für
die Implementierungsdetails bleiben die internen Schrift- und Layout-Subsysteme
zuständig.

\section{Schriften, Schriftsystem-Routing und mehrsprachiger Satz}

\subsection{Schriftfamilien als registrierte Fähigkeiten}

IMPE behandelt einen Schriftnamen nicht lediglich als Ersatz für
\verb|\rmfamily|. Schriftfamilien werden in einem zentralen Katalog mit den
Informationen registriert, die zu ihrer korrekten Verwendung erforderlich
sind: verfügbare Schnitte, Pfade oder Systemschriftnamen,
OpenType-Skript- und Spracheinstellungen, Metadaten für lokale Befehle, globales
Routing-Verhalten sowie, wo nötig, eine spezialisierte Layout- oder
Rendering-Route.

Dieses Modell erlaubt einer einfachen lateinischen Familie und einer Familie
für komplexe Schriftsysteme, dieselbe öffentliche Ladeschnittstelle zu
verwenden, obwohl ihre internen Anforderungen sehr unterschiedlich sein können.
Die gemeinsame Abstraktion ist die \emph{registrierte Familie}, nicht die
Behauptung, jedes Schriftsystem werde identisch implementiert.

\subsection{Automatisches, globales und lokales Laden}

Für den normalen Gebrauch wird eine Familie ohne Modus geladen. IMPE folgt
dann dem registrierten Verhalten der Familie, das lokal, global oder
bereichsspezifisch sein kann. Der Template-Schlüssel \texttt{fonts} wendet
dieselbe automatische Regel auf eine Liste an.

Eine \emph{globale} Familie nimmt an der dokumentweiten Schriftauswahl teil. Bei
einer normalen lateinischen Familie kann dies bedeuten, die Hauptkanäle für
Serifen-, serifenlose und nichtproportionale Schrift zu ersetzen. Bei einer
schriftspezifischen Familie kann es stattdessen bedeuten, nur die Verantwortung
für ausgewählte Unicode-Bereiche zu übernehmen, während Lateinisch, Han und
andere Schriftsysteme unverändert bleiben.

Eine \emph{lokale} Familie wird zur ausdrücklichen Verwendung geladen. Lokale
Auswahl ist nützlich, wenn ein bestimmter Abschnitt eine bestimmte
Schriftsystemschrift, eine regionale Glyphenform oder eine redaktionell
gewählte Schrift benötigt, ohne das dokumentweite Routing zu verändern.

Familien werden beispielsweise so geladen:

\begin{verbatim}
\UseFont{libertinus}
\UseFonts{sanskrit,tibetan,arabic}

\UseFont{libertinus}[global]   % deliberate global override
\UseLocalFonts{sanskrit}       % deliberate local override
\end{verbatim}

oder über die entsprechenden Schlüssel in \verb|\UseTemplateSet|. Die Aliase
\verb|\UseGlobalFont(s)| und \verb|\UseLocalFont(s)| sind explizite
Überschreibungen, nicht die Standardempfehlung. Öffentliche Befehle, die für
einzelne lokale Familien erzeugt werden, sind in deren Katalogregistrierungen
definiert; die vollständige Liste gehört in die Schriftreferenz und nicht in
dieses einführende Handbuch.

Wenn ein expliziter lokaler Schriftbefehl aktiv ist, hat er innerhalb seines
Gültigkeitsbereichs Vorrang vor dem automatischen globalen Bereichs-Router.
Beim Verlassen des lokalen Gültigkeitsbereichs wird der globale Routing-Zustand
wiederhergestellt. Dadurch bleibt eine ausdrückliche Entscheidung des Autors
auch in einem Dokument mit umfangreicher automatischer Schriftsystemauswahl
vorhersehbar.

\subsection{Unicode-Bereichs-Routing}

Einige globale Familien sind bewusst auf bestimmte Bereiche beschränkt. Eine
Devanagari-Familie muss beispielsweise nicht zur lateinischen Hauptschrift
werden, nur weil sie global geladen wird. IMPE kann eine solche Familie nur den
Unicode-Blöcken zuweisen, für die sie zuständig ist. Mehrere registrierte
Familien können daher in einem Dokument nebeneinander bestehen, ohne
fortwährend die Textschrift der jeweils anderen zu ersetzen.

Dieses Routing ist besonders für tatsächlich mehrsprachige Prosa nützlich: Der
Autor schreibt gewöhnlichen Unicode-Text, während IMPE die aktive registrierte
Familie auswählt, sobald der Text in einen konfigurierten Schriftsystembereich
eintritt. Bereichs-Routing hilft bei der Schriftauswahl; das Shaping bleibt
Aufgabe von \XeTeX{}, \texttt{fontspec} und den OpenType-Daten der ausgewählten
Familie.

\subsection{CJK-Behandlung}

CJK-Schriften benötigen eine etwas andere Route als gewöhnliches
Unicode-Bereichs-Umschalten. IMPE verwendet für globale CJK-Familien den
xeCJK-Pfad, sodass die Haupt-, serifenlosen und nichtproportionalen CJK-Kanäle
als CJK-Schriften verwaltet werden können und nicht als Sammlung isolierter
Zeichenumschaltungen.

Gemeinsame Han-Ideogramme bringen ein zusätzliches Problem mit sich: Unicode
allein gibt nicht an, ob ein Han-Zeichen in einem bestimmten Kontext chinesische,
japanische oder koreanische regionale Glyphenkonventionen verwenden soll.
Automatisches Routing weist daher sprachspezifischen Schriftsystemen wie Kana
und Hangul eigene Ansprüche zu, während gemeinsame Han-Zeichen bei der
CJK-Schrift des Dokuments verbleiben. Wenn ausdrücklich eine japanische oder
koreanische Han-Form erforderlich ist, kann ein lokaler Familienbefehl diese
Auswahl direkt treffen.

Diese Unterscheidung ist beabsichtigt. Automatisches Routing soll den Regelfall
bequem machen, ohne vorzugeben, mehrdeutige Han-Zeichen enthielten
Sprachinformationen, die Unicode tatsächlich nicht bereitstellt.

\subsection{Shaping und Zeilenumbruch}

OpenType-Shaping für gewöhnliche komplexe Schriftsysteme wird durch die
registrierten Skript-, Sprach- und Feature-Einstellungen einer Familie
ausgedrückt. Indische, arabische und andere shaping-sensitive Schriften bleiben
daher auf Dokumentebene gewöhnlicher Unicode-Text; IMPE liefert die
Schriftkonfiguration, die die Text-Engine benötigt, um sie korrekt zu formen.

Einige Schriftsysteme benötigen außerdem ein Zeilenumbruchverhalten, das durch
gewöhnliche Leerzeichen nicht ausreichend beschrieben wird. IMPE enthält
schriftsystembewusste Behandlung für Fälle wie Sanskrit in Devanagari,
Tibetisch und Thai. Ziel ist nicht, die Orthographie der jeweiligen Sprache
durch manuelle Umbruchmarkierungen zu ersetzen, sondern den Schrift- und
Verhaltensschichten geeignete Umbruchmöglichkeiten bereitzustellen und dabei
Cluster oder schriftspezifische Interpunktionsregeln zu bewahren.

Rechts-nach-links-Verhalten ist ebenfalls an die Familien gebunden, die es
benötigen. Das öffentliche Schriftmodell bleibt gleich, auch wenn sich das
zugrunde liegende Absatz- oder Inline-Verhalten von einer
links-nach-rechts gesetzten lateinischen Familie unterscheidet.

\subsection{Vertikales und externalisiertes Rendering}

Eine kleine Zahl von Schriftsystemen benötigt mehr als standardmäßiges
horizontales OpenType-Shaping. IMPE enthält daher eine eingebaute vertikale
Layout-Route, die von registrierten vertikalen Familien verwendet wird, und
stellt einen externalisierten Rendering-Pfad für Fälle bereit, in denen ein
Fragment zuverlässiger als kleines zusätzliches \TeX{}-Dokument erzeugt und
anschließend als PDF wieder eingefügt werden kann.

Externalisiertes Rendering wird über einen portablen \texttt{texlua}-Helfer
implementiert. Der Helfer löst die angeforderte \TeX{}-Engine über den
\texttt{PATH} des Systems auf, rendert in ein IMPE-Cacheverzeichnis und gibt
das resultierende PDF an das Hauptdokument zurück. Der Mechanismus bleibt hinter
der Schriftregistrierungsschicht verborgen, sodass ein gewöhnliches Dokument
temporäre Dateien oder Engine-Pfade nicht manuell verwalten muss.

Diese spezialisierten Routen sind Ausnahmen auf Implementierungsebene, nicht auf
der Ebene der Benutzerschnittstelle. Gerade ihre Existenz ermöglicht es, das
öffentliche Schriftmodell kohärent zu halten, wenn ein Schriftsystem nicht mit
einer einfachen \verb|\newfontfamily|-Deklaration behandelt werden kann.

\subsection{Wo die vollständige Schriftreferenz zu finden ist}

Der detaillierte Familienkatalog, die Registrierungssyntax, Regeln für
Stil-Fallbacks, das globale Routing-Modell, vertikale Parameter, das
externalisierte Backend und schriftspezifische Hinweise werden in
\texttt{docs/FONTS.md} gepflegt. Benutzer, die eine neue Familie hinzufügen
oder das Font-Routing untersuchen, sollten dieses Dokument als technische
Referenz behandeln.

\section{Layouts und Features}

\subsection{Layout-Voreinstellungen}

Layouts beschreiben die Darstellung auf Dokumentebene. Eine öffentliche
Layout-Voreinstellung ist ein strukturiertes Bündel aus Seitengeometrie,
Textabständen, Kopfzeilenverhalten, Buchverhalten oder Folienverhalten. Die
einzelnen Low-Level-Komponenten sind intern; gewöhnliche Dokumente wählen die
öffentliche Voreinstellung als Ganzes.

Die öffentlichen Voreinstellungen in IMPE 1.0.0 sind:

\begin{description}
  \item[\texttt{en\_doc}] Englisches Artikel-/Report-Layout.
  \item[\texttt{zh\_doc}] Chinesisches Artikel-/Report-Layout.
  \item[\texttt{en\_book}] Englisches Buch-/Report-Layout mit buchorientiertem
    Seiten- und Kolumnentitelverhalten.
  \item[\texttt{zh\_book}] Chinesisches Gegenstück für Buch/Report.
  \item[\texttt{beamer}] Präsentationslayout für Beamer-Dokumente.
\end{description}

Voreinstellungen deklarieren, welche zugrunde liegenden Dokumentklassen sie
unterstützen. IMPE prüft diese Kompatibilität, bevor die Voreinstellung
angewendet wird, sodass ein Folienlayout nicht stillschweigend auf einen
Artikel angewendet wird oder umgekehrt.

Wrapper-Klassen wählen das entsprechende Layout automatisch. Eine explizite
Auswahl ist bei rohen Standardklassen oder dann nützlich, wenn ein Report eine
buchorientierte Voreinstellung verwenden soll:

\begin{verbatim}
\UseLayout{en_book}
\end{verbatim}

Die genauen Komponentenwerte und die Geometrie sind in
\texttt{docs/LAYOUTS.md} dokumentiert.

\subsection{Feature-Module}

Features sind flache, kombinierbare Module. Anders als Layouts besitzen sie
keine zweite Voreinstellungsebene. Eine Feature-ID lädt beim ersten Anfordern
einfach die entsprechende optionale Funktionalität; wiederholte Anforderungen
werden ignoriert.

Die öffentliche Feature-Menge besteht derzeit aus:

\begin{description}
  \item[\texttt{math}] Standardpakete für Mathematik, Theorem-Umgebungen und
    optionale Auswahl der Mathematikschrift.
  \item[\texttt{hyperlinks}] Hyperlinks, PDF-Lesezeichen, stabile Ziele und
    Repository-Hyperlink-Standards.
  \item[\texttt{citations}] Unterstützung für Zitate mit
    \texttt{biblatex}/\texttt{biber} und auswählbaren
    Zitationsstil-Voreinstellungen.
  \item[\texttt{index}] Registererzeugung und IMPE-Helfer für Begriffe und
    Register.
  \item[\texttt{tables}] Tabellenpakete, wiederverwendbare Spaltentypen und
    Tabellenumgebungen.
  \item[\texttt{image}] Helfer für Abbildungen, Bilder und Mehrfachpanels.
  \item[\texttt{lists\_envs}] wiederverwendbare Formatierung für Beispielblöcke.
  \item[\texttt{headers}] konfigurierbare Kolumnentitel für artikel-, report-
    und buchartige Dokumente.
\end{description}

Ein Dokument lädt jede benötigte Kombination:

\begin{verbatim}
\UseFeatures{math,hyperlinks,tables,image}
\end{verbatim}

oder

\begin{verbatim}
\UseTemplateSet{
  features = {citations,index,hyperlinks}
}
\end{verbatim}

\subsection{Mathematik}

Das Feature \texttt{math} lädt den standardmäßigen AMS-orientierten
Mathematik-Stack und definiert theoremartige Umgebungen wie \texttt{theorem},
\texttt{lemma}, \texttt{proposition}, \texttt{corollary},
\texttt{definition}, \texttt{example} und \texttt{remark}. Das Laden einer
Textschrift erzwingt nicht von selbst ein neues mathematisches Alphabet:
Computer-Modern-Mathematik bleibt Standard, sofern nicht ausdrücklich eine
andere Mathematikroute gewählt wird.

Eine andere Mathematikfamilie kann vor dem Laden des Features ausgewählt
werden. Zum Beispiel

\begin{verbatim}
\UseMathFont{libertinus}
\UseFeature{math}
\end{verbatim}

verwendet Libertinus Math über \texttt{unicode-math}. Weitere unterstützte
Routen und der genaue Paket-Stack sind in \texttt{docs/FEATURES.md}
aufgeführt.

\subsection{Zitate, Register, Tabellen und Abbildungen}

Die Feature-Module sollen nützliche Standards bereitstellen, ohne das
zugrunde liegende \LaTeX{}-Ökosystem zu verbergen. Das Feature
\texttt{citations} verwendet beispielsweise \texttt{biblatex} und
\texttt{biber}; Dokumente verwenden weiterhin gewöhnliche Befehle wie
\verb|\addbibresource|, \verb|\textcite| und \verb|\printbibliography|.
Zitations-Voreinstellungen wie APA, GB/T 7714, numerisch oder Autor--Jahr können
vor dem Laden des Features gewählt werden.

Das Feature \texttt{index} baut auf \texttt{imakeidx} auf und ergänzt den
Helfer \verb|\Term|, um einen Begriff zugleich zu setzen und zu indizieren.
Die Features \texttt{tables} und \texttt{image} stellen auf ähnliche Weise
praktische höherstufige Umgebungen bereit, während native Befehle von
\texttt{booktabs}, \texttt{graphicx} und verwandten Paketen verfügbar bleiben.

Dies ist im gesamten Feature-System die beabsichtigte Balance: IMPE soll
wiederholte Einrichtung reduzieren und keine private Ersatzsprache für
gewöhnliches \LaTeX{} schaffen.

\subsection{Kopfzeilen und Hyperlinks}

Das Feature \texttt{headers} stellt über \texttt{fancyhdr} Kolumnentitel
bereit. Standardmäßig leitet es aus der ersten Zeile von \verb|\title| einen
festen Kurztitel ab; \verb|\HeaderTitle{...}| kann verwendet werden, wenn eine
kürzere Form bevorzugt wird. Artikelartige Dokumente verwenden
Abschnittsinformationen für den wechselnden Kolumnentitel, report- und
buchartige Dokumente Kapitelinformationen.

Das Feature \texttt{hyperlinks} lädt \texttt{hyperref} und
\texttt{bookmark} mit Unicode-tauglichen PDF-Einstellungen, ausgeblendeten
Linkrahmen, verlinkten Inhaltsverzeichnisstrukturen und stabilen Zielen.
Dokumente können weiterhin gewöhnliche Befehle wie \verb|\label|, \verb|\ref|,
\verb|\href| und \verb|\url| verwenden.

Für die vollständige Liste der feature-spezifischen Befehle und Umgebungen
sollte \texttt{docs/FEATURES.md} konsultiert werden; dieser Abschnitt ist kein
API-Katalog.

\section{Kompatibilität und Migration}\label{sec:compatibility}

IMPE 1.0.0 etabliert die \texttt{impe*}-Namen als kanonische öffentliche
Schnittstelle. Neue Dokumente sollten daher \texttt{impe.sty},
\texttt{impeart}, \texttt{impebook}, \texttt{impereport},
\texttt{impebeamer} und ihre chinesischen Varianten verwenden.

Frühere Entwicklungsversionen verwendeten das Namensschema \texttt{next*}.
Diese Einstiegspunkte bleiben in 1.0.0 als Kompatibilitäts-Wrapper unterstützt;
sie sind in dieser Version nicht als veraltet markiert. Bestehende Dokumente
wie

\begin{verbatim}
\documentclass{nextart}
\usepackage{nextsystem}
\end{verbatim}

können daher weiterhin gegen dieselbe zugrunde liegende Implementierung
kompiliert werden. Die Kompatibilitätsschicht dient der Bewahrung alter
Dokumente und nicht der Definition eines parallelen API-Satzes für neue
Arbeiten.

Der direkte Installer erkennt außerdem eine verwaltete Installation der
0.1.x-Reihe unter \texttt{tex/latex/nextsystem/}. Bei der Migration installiert
er die kanonische Laufzeitumgebung unter \texttt{tex/latex/impe/}, migriert
erkannte lokale Override-Dateien, wo dies angemessen ist, entfernt Dateien, die
bekanntermaßen zur alten verwalteten Laufzeit gehören, und lässt nicht
verwandte benutzererstellte Dateien oder Unterverzeichnisse unverändert.

Eine alte \texttt{nextsystem.local.tex} kann die Migration daher überstehen,
ohne dass die alten generischen Runtime-Dateinamen im TEXMF-Baum sichtbar
bleiben müssen. Neue lokale Konfiguration sollte \texttt{impe.local.tex}
verwenden.

Eine Konvertierung des Quelltexts ist nicht erforderlich, nur weil ein Dokument
mit den alten Wrapper-Namen geschrieben wurde. Autoren können diese Namen bei
Gelegenheit aktualisieren; Kompatibilität ist bewusst eine Angelegenheit der
Laufzeitumgebung und keine erzwungene Migration jedes bestehenden Dokuments.

\section{Paketstruktur und Erweiterungspunkte}

Das Repository trennt öffentliche Einstiegspunkte, stabile Mechanismen,
Registrierungen, spezialisierte Implementierungen und lokale Ressourcen:

\begin{verbatim}
package/   public package and class entry points
core/      stable subsystem mechanisms
catalog/   public ids and registration data
modules/   extendable or specialized implementations
assets/    local runtime resources such as fonts
scripts/   installation and release tooling
doc/       manual sources
docs/      detailed subsystem documentation
\end{verbatim}

\subsection{Der stabile Kern}

Die Schicht \texttt{core/} enthält das Framework, das nicht geändert werden
muss, wenn eine neue Schriftfamilie, Layout-Voreinstellung oder ein neues
Feature hinzugefügt wird. Sie implementiert die Schriftregistrierung und
Routing-Mechanik, die Layout-Registrierung und Kompatibilitätsprüfungen, den
Feature-Loader und damit verbundenes gemeinsames Verhalten.

Diese Unterscheidung ist für die Wartung wichtig. Wenn das Hinzufügen einer
gewöhnlichen OpenType-Schrift eine Änderung an der Font-Routing-Engine erfordert
oder das Hinzufügen eines Features den generischen Feature-Loader verändern
muss, ist das Erweiterungsmodell gescheitert. Der Kern ist daher bewusst vom
wachsenden Katalog konkreter Ressourcen getrennt.

\subsection{Kataloge}

Die Schicht \texttt{catalog/} vergibt öffentliche IDs und Metadaten. Sie enthält
die zentralen Registrierungen für Schriften, Layouts und Features. Für viele
Erweiterungen ist dies der richtige Ausgangspunkt.

Eine gewöhnliche Schriftfamilie, die nur standardmäßiges
\texttt{fontspec}-Shaping benötigt, kann in der Regel vollständig als
Katalogregistrierung ausgedrückt werden. Ebenso wird eine öffentliche
Layout-Voreinstellung aus vorhandenen internen Komponenten zusammengesetzt,
und eine Feature-ID verweist auf ihr Implementierungsmodul. Die zentrale
Ablage dieser Deklarationen macht die öffentliche Oberfläche überprüfbar, ohne
den gesamten Quellbaum durchsuchen zu müssen.

\subsection{Module}

Die Schicht \texttt{modules/} ist Funktionalität vorbehalten, die tatsächlich
eine eigene Implementierung benötigt: Feature-Module, spezialisierte
Schriftsystemunterstützung oder wiederverwendbare interne Layout-Komponenten.
Ein neues Modul sollte daher eine echte Verhaltenserweiterung darstellen und
nicht lediglich ein bequemer Ort für einen Schriftnamen sein.

Diese Regel ist für mehrsprachige Unterstützung besonders wichtig. Die meisten
Schriftsysteme sollten gewöhnliches OpenType-Shaping und Katalogmetadaten
verwenden. Spezialisierte Module sind nur dann angebracht, wenn das
Schriftsystem oder die Rendering-Strategie ein Verhalten erfordert, das der
generische Kern nicht sauber ausdrücken kann.

\subsection{IMPE erweitern}

Die detaillierten Registrierungsschnittstellen liegen bewusst außerhalb des
Rahmens dieses Benutzerhandbuchs. Entwickler, die das System erweitern möchten,
sollten mit \texttt{docs/SYSTEM.md} beginnen und anschließend das relevante
Subsystemdokument konsultieren: \texttt{docs/FONTS.md},
\texttt{docs/LAYOUTS.md} oder \texttt{docs/FEATURES.md}.

Die allgemeine Regel ist einfach: Eine Registrierung ist einem neuen Mechanismus
vorzuziehen, und bei tatsächlich außergewöhnlichem Verhalten ist ein
spezialisiertes Modul einer Änderung des stabilen Kerns vorzuziehen.

\section{Weiterführende Dokumentation und Rückmeldungen}

Dieses Handbuch ist der normale Ausgangspunkt für die Verwendung von IMPE. Das
Repository enthält detailliertere technische Dokumentation für Benutzer, die
eine vollständige Referenz benötigen oder das System erweitern möchten.

\begin{description}
  \item[\texttt{docs/SYSTEM.md}] Architektur, öffentliche Einstiegsschichten,
    Template-Konfiguration, Release-Modell und Erweiterungsgrenzen.
  \item[\texttt{docs/FONTS.md}] Familienregistrierung, lokale und globale Modi,
    Routing, Fallback, Shaping, Schreibverhalten, vertikale Unterstützung und
    externalisiertes Rendering.
  \item[\texttt{docs/LAYOUTS.md}] öffentliche Layout-Voreinstellungen,
    Kompatibilitätsziele und das Komponentenmodell, aus dem Voreinstellungen
    aufgebaut werden.
  \item[\texttt{docs/FEATURES.md}] öffentliche Feature-IDs sowie Befehle,
    Umgebungen, Paket-Stacks und Kompilierungsanforderungen der einzelnen
    Module.
  \item[\texttt{CHANGELOG.md}] Änderungen in veröffentlichten Versionen.
  \item[Showcase] der vollständige visuelle Showcase ist im
    \href{\IMPEShowcaseURL}{Projekt-Repository} verfügbar.
\end{description}

Das Projekt-Repository ist erreichbar unter

\begin{center}
\href{\IMPERepo}{\texttt{github.com/KelvinYangBK67/IMPE-LaTeX-System}}
\end{center}

Fehlerberichte, Feature-Wünsche, Dokumentationskorrekturen,
Portabilitätsberichte und andere konkrete Rückmeldungen sind über GitHub Issues
willkommen:

\begin{center}
\href{\IMPEIssues}{\texttt{github.com/KelvinYangBK67/IMPE-LaTeX-System/issues}}
\end{center}

Bei der Meldung eines Satzproblems ist es hilfreich, eine minimale Quelldatei,
die \TeX{}-Engine und \TeX{}-Live-Version, die IMPE-Version und die relevanten
Schriftinformationen anzugeben. Berichte zu einer lokalen oder privat
bereitgestellten Schrift sollten, soweit möglich, außerdem zwischen einem
Problem der Glyphenabdeckung und einem IMPE-Routing-Problem unterscheiden.

Das englische Handbuch ist der maßgebliche Handbuchtext. Ausgaben in anderen
Sprachen dürfen Prosa und Terminologie an die jeweiligen technischen
Konventionen anpassen, doch ihre Befehle, Beispiele und die Beschreibung des
IMPE-Verhaltens sollen mit dieser Ausgabe synchron bleiben.

\appendix

\section{Kurzreferenz}
\label{app:quick-reference}

Dieser Anhang fasst die dokumentierte öffentliche Schnittstelle von IMPE 1.0.0
zusammen. Er dient zum Nachschlagen und nicht als Ersatz für die Erläuterungen
im Haupttext oder die detaillierte Subsystemdokumentation unter
\texttt{docs/}. Interne Implementierungsbefehle, insbesondere Befehle, die mit
\texttt{\string\Next...} beginnen, sind nicht Teil dieser Referenz.

\subsection{Empfohlenes Modell zum Laden von Schriften}

Für den normalen Gebrauch sollte innerhalb von
\texttt{\string\UseTemplateSet} vorzugsweise
\texttt{fonts = \{...\}} oder direkt \texttt{\string\UseFont\{...\}} verwendet
werden. Ohne expliziten Modus folgt IMPE dem registrierten Verhalten der
Familie: Je nach Bedarf kann ein lokaler Befehl definiert, eine dokumentweite
Schrift aktiviert oder ein skript- bzw. bereichsspezifisches automatisches
Routing eingerichtet werden.

Ein expliziter Modus sollte nur verwendet werden, wenn dieses automatische
Verhalten nicht den Anforderungen des Dokuments entspricht:

\begin{verbatim}
\UseFont{sanskrit}          % recommended: registered automatic behavior
\UseFont{sanskrit}[local]   % explicitly local only
\UseFont{sanskrit}[global]  % explicitly global/range-global

\UseTemplateSet{
  fonts = {sanskrit,tibetan,arabic},
  features = {math,hyperlinks}
}
\end{verbatim}

Die Komfortbefehle \texttt{\string\UseGlobalFont(s)} und
\texttt{\string\UseLocalFont(s)} drücken dieselbe explizite Unterscheidung aus.
Die Template-Schlüssel \texttt{globalfonts} und \texttt{mainfonts} sind
ebenfalls ausdrückliche globale Überschreibungen; sie ersetzen im Normalfall
nicht \texttt{fonts}.

\subsection{Dokumentklassen und Paket-Einstiegspunkte}

\begin{longtable}{@{}p{0.31\linewidth}p{0.63\linewidth}@{}}
\toprule
\textbf{Einstiegspunkt} & \textbf{Zweck} \\
\midrule
\endfirsthead
\toprule
\textbf{Einstiegspunkt} & \textbf{Zweck} \\
\midrule
\endhead
\texttt{impeart} &
Englischer Artikel-Wrapper; verwendet standardmäßig \texttt{en\_doc} mit CMU als globaler Basisfamilie. \\
\texttt{impeart\_zh} &
Chinesischer Artikel-Wrapper; verwendet standardmäßig \texttt{zh\_doc} mit CMU und Shanggu. \\
\texttt{impereport} &
Englischer Report-Wrapper; verwendet standardmäßig \texttt{en\_doc}. \\
\texttt{impereport\_zh} &
Chinesischer Report-Wrapper; verwendet standardmäßig \texttt{zh\_doc}. \\
\texttt{impebook} &
Englischer Buch-Wrapper; verwendet standardmäßig \texttt{en\_book}. \\
\texttt{impebook\_zh} &
Chinesischer Buch-Wrapper; verwendet standardmäßig \texttt{zh\_book}. \\
\texttt{impebeamer} &
Englischer Beamer-Wrapper; verwendet standardmäßig das Layout \texttt{beamer}. \\
\texttt{impebeamer\_zh} &
Chinesischer Beamer-Wrapper; verwendet standardmäßig das Layout \texttt{beamer} mit CMU und Shanggu. \\
\texttt{impe.sty} &
Lädt IMPE unter einer gewöhnlichen \LaTeX{}-Klasse, wenn keine Wrapper-Klasse gewünscht ist. \\
\bottomrule
\end{longtable}

Die alten \texttt{next*}-Paket- und Klassennamen bleiben in Version 1.0.0
Kompatibilitätsaliase; neue Dokumente sollten jedoch die kanonischen
\texttt{impe*}-Einstiegspunkte verwenden.

\subsection{Zentrale Konfigurationsbefehle}

\begin{longtable}{@{}p{0.38\linewidth}p{0.56\linewidth}@{}}
\toprule
\textbf{Befehl} & \textbf{Zweck} \\
\midrule
\endfirsthead
\toprule
\textbf{Befehl} & \textbf{Zweck} \\
\midrule
\endhead
\texttt{\string\UseTemplateSet\{...\}} &
Konfiguriert Layout, Schriften und Features in einer Deklaration. \\
\texttt{\string\UseFont\{id\}[mode]} &
Lädt eine registrierte Schriftfamilie; \texttt{mode} wird für das empfohlene automatische Verhalten weggelassen oder explizit als \texttt{local}/\texttt{global} angegeben. \\
\texttt{\string\UseFonts\{ids\}[mode]} &
Lädt mehrere registrierte Schriftfamilien mit demselben optionalen Modus. \\
\texttt{\string\UseGlobalFont\{id\}} &
Fordert ausdrücklich den globalen Modus einer Familie an. \\
\texttt{\string\UseGlobalFonts\{ids\}} &
Fordert ausdrücklich den globalen Modus mehrerer Familien an. \\
\texttt{\string\UseMainFont\{id\}} &
Alias von \texttt{\string\UseGlobalFont}; als dokumentorientierter Name beibehalten. \\
\texttt{\string\UseMainFonts\{ids\}} &
Alias von \texttt{\string\UseGlobalFonts}. \\
\texttt{\string\UseLocalFont\{id\}} &
Fordert ausdrücklich den lokalen Modus einer Familie an. \\
\texttt{\string\UseLocalFonts\{ids\}} &
Fordert ausdrücklich den lokalen Modus mehrerer Familien an. \\
\texttt{\string\UseLayout\{id\}} &
Lädt eine registrierte Layout-Voreinstellung. \\
\texttt{\string\UseLayouts\{ids\}} &
Lädt mehrere registrierte Layout-Voreinstellungen. \\
\texttt{\string\UseFeature\{id\}} &
Lädt ein optionales Feature-Modul. \\
\texttt{\string\UseFeatures\{ids\}} &
Lädt mehrere optionale Feature-Module. \\
\texttt{\string\SetCatalogFontRoot\{path\}} &
Überschreibt das Stammverzeichnis, aus dem katalogisierte Schriftdateien aufgelöst werden. \\
\texttt{\string\subtitle\{text\}} &
Setzt den Dokumentuntertitel, der vom IMPE-Titelblock verwendet wird. \\
\texttt{\string\HeaderTitle\{text\}} &
Überschreibt den festen Titel, den das Feature \texttt{headers} verwendet. \\
\texttt{\string\HeaderStyle\{style\}} &
Wählt den Kopfzeilenstil; dokumentierte Werte sind \texttt{running} und \texttt{title}. \\
\bottomrule
\end{longtable}

\subsection{Schlüssel des Template-Sets}

\begin{longtable}{@{}p{0.23\linewidth}p{0.69\linewidth}@{}}
\toprule
\textbf{Schlüssel} & \textbf{Bedeutung} \\
\midrule
\endfirsthead
\toprule
\textbf{Schlüssel} & \textbf{Bedeutung} \\
\midrule
\endhead
\texttt{layout} &
Wählt eine öffentliche Layout-Voreinstellung. \\
\texttt{fonts} &
Lädt eine oder mehrere registrierte Familien mit ihrem automatischen registrierten Verhalten; dies ist der empfohlene Schlüssel für normales Schriftladen. \\
\texttt{globalfonts} &
Zwingt die aufgeführten Familien ausdrücklich durch ihren globalen Modus oder globalen Routing-Pfad. \\
\texttt{mainfonts} &
Alias von \texttt{globalfonts}. \\
\texttt{features} &
Lädt ein oder mehrere optionale Feature-Module. \\
\bottomrule
\end{longtable}

\subsection{Öffentliche Layout-IDs}

\begin{longtable}{@{}p{0.23\linewidth}p{0.69\linewidth}@{}}
\toprule
\textbf{ID} & \textbf{Zweck} \\
\midrule
\endfirsthead
\toprule
\textbf{ID} & \textbf{Zweck} \\
\midrule
\endhead
\texttt{en\_doc} &
Englisches Artikel-/Report-Layout mit A4-Geometrie und englischen Textabständen. \\
\texttt{zh\_doc} &
Chinesisches Artikel-/Report-Layout mit A4-Geometrie und chinesischen Textabständen. \\
\texttt{en\_book} &
Englisches Buch-/Report-Layout mit Buchgeometrie, Kolumnentiteln und Kapitelbeginn auf rechten Seiten. \\
\texttt{zh\_book} &
Chinesisches Buch-/Report-Layout mit Buchgeometrie, chinesischen Textabständen, Kolumnentiteln und Kapitelbeginn auf rechten Seiten. \\
\texttt{beamer} &
Präsentationslayout mit IMPE-Beamer-Typographie und Navigationseinstellungen. \\
\bottomrule
\end{longtable}

\subsection{Öffentliche Feature-IDs}

\begin{longtable}{@{}p{0.23\linewidth}p{0.69\linewidth}@{}}
\toprule
\textbf{ID} & \textbf{Zweck} \\
\midrule
\endfirsthead
\toprule
\textbf{ID} & \textbf{Zweck} \\
\midrule
\endhead
\texttt{math} &
Lädt den IMPE-Mathematik-Stack und Theorem-Umgebungen. \\
\texttt{hyperlinks} &
Lädt Hyperlinks, PDF-Lesezeichen, verlinkte Überschriften und bidirektionale Fußnotenlinks. \\
\texttt{citations} &
Lädt Zitationsunterstützung mit \texttt{biblatex}/\texttt{csquotes} und IMPE-Helfer für Zitationsstile. \\
\texttt{index} &
Lädt Registerunterstützung und den Helfer \texttt{\string\Term}. \\
\texttt{tables} &
Lädt den Tabellen-Stack, Spaltentypen und IMPE-Tabellenumgebungen. \\
\texttt{image} &
Lädt Helfer für Bilder, Bildunterschriften und Panel-Abbildungen. \\
\texttt{lists\_envs} &
Lädt die Anzeigeumgebung \texttt{ExampleBlock}. \\
\texttt{headers} &
Lädt konfigurierbare Kolumnentitel für artikel-, report- und buchartige Dokumente. \\
\bottomrule
\end{longtable}

\subsection{Feature-spezifische Befehle und Werte}

\begin{longtable}{@{}p{0.51\linewidth}p{0.43\linewidth}@{}}
\toprule
\textbf{Befehl oder Einstellung} & \textbf{Zweck / dokumentierte Werte} \\
\midrule
\endfirsthead
\toprule
\textbf{Befehl oder Einstellung} & \textbf{Zweck / dokumentierte Werte} \\
\midrule
\endhead
\texttt{\string\UseMathFont\{value\}} &
Wählt vor dem Laden von \texttt{math} die Mathematikschrift-Route: \texttt{auto}, \texttt{libertinus}, \texttt{newcm}, \texttt{mlmodern} oder einen anderen Schriftnamen, der an \texttt{\string\setmathfont} übergeben wird. \\
\texttt{\string\UseCitationStyle\{value\}} &
Wählt vor dem Laden von \texttt{citations} eine Zitationsvoreinstellung: \texttt{APA}, \texttt{GB}, \texttt{numeric} oder \texttt{author-year}. \\
\texttt{\string\SetCitationBiblatexOptions\{...\}} &
Ersetzt vor dem Laden von \texttt{citations} die wirksame \texttt{biblatex}-Optionsliste. \\
\texttt{\string\HeaderTitle\{text\}} &
Setzt einen kurzen festen Titel für Kolumnentitel. \\
\texttt{\string\HeaderStyle\{value\}} &
Wählt \texttt{running} oder \texttt{title}. \\
\texttt{\string\IndexTitle} &
Optionale Überschreibung des Registertitels, die vor dem Laden von \texttt{index} definiert wird. \\
\texttt{\string\Term[options]\{display\}[description]} &
Setzt einen Begriff fett und indiziert sein erstes Auftreten; Optionen umfassen \texttt{sort}, \texttt{key} und \texttt{parentheses}. \\
\texttt{\string\TablesSetup} &
Wendet die IMPE-Tabellenkonfiguration an oder erneut an. \\
\texttt{\string\Panel} &
Fügt ein Panel innerhalb von \texttt{PanelFigure} oder \texttt{PanelFigure*} ein. \\
\texttt{\string\TemplateFigurePaths} &
Konfiguriert die Standardsuchpfade des Image-Features. \\
\texttt{\string\OneImageDefaultWidth} &
Setzt die Standardbreite für \texttt{OneImage}. \\
\texttt{\string\OneImageMaxHeight} &
Setzt die maximale Standardbildhöhe für \texttt{OneImage}. \\
\texttt{\string\OneImageDefaultPlacement} &
Setzt die Standard-Float-Platzierung für \texttt{OneImage}. \\
\texttt{\string\PanelDefaultCols} &
Setzt die Standardzahl der Spalten für Panel-Abbildungen. \\
\texttt{\string\PanelDefaultHeight} &
Setzt die Standardhöhe eines Panels. \\
\texttt{\string\PanelDefaultMode} &
Setzt den Standardmodus des Panel-Layouts. \\
\texttt{\string\PanelDefaultPlacement} &
Setzt die Standard-Float-Platzierung für Panel-Abbildungen. \\
\bottomrule
\end{longtable}

Standardbefehle, die unmittelbar von Paketen bereitgestellt werden, welche ein
Feature lädt --- etwa \texttt{\string\addbibresource},
\texttt{\string\printbibliography}, \texttt{\string\printindex} oder die
Regelbefehle von \texttt{booktabs} --- behalten ihre normale Paketsyntax und
werden von IMPE nicht neu definiert.

\subsection{Öffentliche Umgebungen und Tabellenspaltentypen}

\begin{longtable}{@{}p{0.37\linewidth}p{0.57\linewidth}@{}}
\toprule
\textbf{Name} & \textbf{Zweck} \\
\midrule
\endfirsthead
\toprule
\textbf{Name} & \textbf{Zweck} \\
\midrule
\endhead
\texttt{theorem}, \texttt{lemma}, \texttt{proposition}, \texttt{corollary} &
Nummerierte theoremartige Umgebungen des Features \texttt{math}. \\
\texttt{definition}, \texttt{example}, \texttt{remark} &
Weitere Theorem-Umgebungen des Features \texttt{math}. \\
\texttt{TableInlineFit} &
Kompakte Tabellenumgebung, die eine Inline-Tabelle an die verfügbare Breite anpassen soll. \\
\texttt{TableLong} &
Longtable-Umgebung für Material, das sich über mehrere Seiten erstrecken kann. \\
\texttt{TableBook}, \texttt{TableBookX} &
Tabellenumgebungen im Buchstil. \\
\texttt{TableBookNotes} &
Tabellenumgebung im Buchstil mit Anmerkungen. \\
\texttt{NiceBooktable}, \texttt{NiceBooktableX}, \texttt{NiceBooktableNotes} &
Alternative Tabellenhelfer im Buchstil des Features \texttt{tables}. \\
\texttt{L}, \texttt{C}, \texttt{R} &
Linksbündige, zentrierte und rechtsbündige \texttt{tabularx}-Spalten. \\
\texttt{P\{w\}}, \texttt{M\{w\}}, \texttt{B\{w\}} &
Absatzspalten fester Breite mit linksbündiger, zentrierter und rechtsbündiger Ausrichtung. \\
\texttt{OneImage} &
Einzelbild-Abbildung; unter Beamer inline statt als Float. \\
\texttt{OneImageInline} &
Zentriertes, nicht schwebendes Bild. \\
\texttt{PanelFigure} &
Beschriftete Mehrfachpanel-Abbildung mit Unterbeschriftungen außerhalb von Beamer. \\
\texttt{PanelFigure*} &
Unbeschriftete Mehrfachpanel-Abbildung. \\
\texttt{ExampleBlock} &
Eingerückter kursiver Anzeige-Block für Beispiele, Zitate, linguistische Daten oder Lehrmaterial. \\
\bottomrule
\end{longtable}

\subsection{Registrierte Schriftfamilien-IDs}

Die folgende Tabelle listet die von \texttt{\string\UseFont} in IMPE 1.0.0
akzeptierten Schriftfamilien-IDs auf. Die Spalte \textbf{Auto} fasst zusammen,
was ein Aufruf von \texttt{\string\UseFont\{id\}} ohne Modus normalerweise
aktiviert:

\begin{description}
  \item[\texttt{G}] dokumentweite globale Familie;
  \item[\texttt{L}] lokale Familie/lokaler Befehl;
  \item[\texttt{L+G}] lokaler Zugriff plus vollständige dokumentweite globale Aktivierung;
  \item[\texttt{L+R}] lokaler Zugriff plus automatisches skript- oder bereichsbeschränktes globales Routing.
\end{description}

Die automatische Route ist die empfohlene Schnittstelle. Eine ausdrückliche
Anforderung \texttt{[local]} oder \texttt{[global]} umgeht bewusst einen Teil
dieser automatischen Entscheidung und sollte daher nur verwendet werden, wenn
sie erforderlich ist.

\begin{longtable}{@{}p{0.30\linewidth}p{0.11\linewidth}p{0.18\linewidth}p{0.31\linewidth}@{}}
\toprule
\textbf{Familien-ID} & \textbf{Auto} & \textbf{Lokaler Befehl} & \textbf{Typische Verwendung} \\
\midrule
\endfirsthead
\toprule
\textbf{Familien-ID} & \textbf{Auto} & \textbf{Lokaler Befehl} & \textbf{Typische Verwendung} \\
\midrule
\endhead
\texttt{cmu} & G & --- & Computer-Modern-Unicode-Basisfamilie. \\
\texttt{noto} & L+G & \texttt{\string\NOT} & Allgemeine Noto-Lateinfamilie. \\
\texttt{times} & L+G & \texttt{\string\TIM} & Systemfamilien-Route für Times/Arial/Consolas. \\
\texttt{gentium} & L+G & \texttt{\string\GEN} & Gentium Plus. \\
\texttt{charis} & L+G & \texttt{\string\CHA} & Charis SIL. \\
\texttt{libertinus} & L+G & \texttt{\string\LIB} & Libertinus-Familie für Serif/Sans/Mono. \\
\texttt{mlmodern} & L & \texttt{\string\MLM} & Lokale Route für Latin Modern / MLModern. \\

\texttt{anatolian} & L & \texttt{\string\CA} & Karisches / anatolisches Schriftmaterial. \\
\texttt{coptic} & L & \texttt{\string\CO} & Koptisch. \\
\texttt{bopomofo} & L & \texttt{\string\ZY} & Bopomofo. \\
\texttt{cuneiform} & L & \texttt{\string\CU} & Keilschrift. \\
\texttt{glagolitic} & L & \texttt{\string\GL} & Glagolitisch. \\
\texttt{italic} & L & \texttt{\string\OI} & Altitalisch. \\
\texttt{hungarian} & L & \texttt{\string\OH} & Altungarisch. \\
\texttt{runic} & L & \texttt{\string\RU} & Runenschrift. \\
\texttt{armenian} & L & \texttt{\string\HY} & Armenisch. \\
\texttt{georgian} & L & \texttt{\string\KA} & Georgisch. \\

\texttt{hindi} & L+R & \texttt{\string\HI} & Hindi-Devanagari mit registriertem Bereichs-Routing. \\
\texttt{sanskrit} & L+R & \texttt{\string\SA} & Sanskrit-Devanagari mit Sanskrit-bewusstem Routing. \\
\texttt{devanagari} & L & \texttt{\string\DEV} & Generisches Devanagari ohne sprachspezifisches Verhalten. \\
\texttt{tamil} & L & \texttt{\string\TA} & Tamil. \\
\texttt{brahmi} & L & \texttt{\string\BR} & Brahmi. \\
\texttt{tibetan} & L+R & \texttt{\string\TI} & Tibetisch mit Tsheg-bewusstem Zeilenumbruch. \\
\texttt{segoe} & L & \texttt{\string\SEG} & Lokale Route für Segoe Historic. \\
\texttt{thai} & L & \texttt{\string\TH} & Thai. \\

\texttt{arabic} & L & \texttt{\string\AR} & Arabisch; Serif/Sans und alternative Stilvarianten. \\
\texttt{urdu} & L & \texttt{\string\UR} & Urdu-Nastaliq. \\
\texttt{aramaic} & L & \texttt{\string\IA} & Reichsaramäisch. \\
\texttt{nabataean} & L & \texttt{\string\NB} & Nabatäisch. \\
\texttt{hebrew} & L & \texttt{\string\HE} & Hebräisch. \\
\texttt{syriac} & L & \texttt{\string\SY} & Syrisch. \\
\texttt{syriac\_eastern} & L+R & \texttt{\string\SYE} & Ostsyrisch mit registriertem Bereichs-Routing. \\
\texttt{kharosthi} & L & \texttt{\string\KH} & Kharoshthi. \\
\texttt{khitan\_small} & L & \texttt{\string\KHS} & Kleine Khitan-Schrift mit ihrem spezialisierten Kompositor. \\
\texttt{pahlavi\_parthian} & L+R & \texttt{\string\PAR} & Inschriftliches Parthisch. \\
\texttt{pahlavi\_inscriptional} & L+R & \texttt{\string\PAH} & Inschriftliches Pahlavi. \\
\texttt{pahlavi\_psalter} & L & \texttt{\string\PSP} & Psalter-Pahlavi mit spezialisierter Unterstützung. \\
\texttt{avestan} & L & \texttt{\string\AV} & Avestisch. \\
\texttt{manichaean} & L & \texttt{\string\MA} & Manichäisch. \\
\texttt{phoenician} & L & \texttt{\string\PH} & Phönizisch. \\
\texttt{samaritan} & L & \texttt{\string\SM} & Samaritanisch. \\
\texttt{sogdian} & L & \texttt{\string\SG} & Sogdisch. \\
\texttt{sogdian\_old} & L+R & \texttt{\string\SGO} & Altsogdisch mit registriertem Bereichs-Routing. \\

\texttt{chinese\_simplified} & L+R & \texttt{\string\SC} & CJK-Familie für vereinfachtes Chinesisch. \\
\texttt{chinese\_traditional} & L+R & \texttt{\string\TC} & CJK-Familie für traditionelles Chinesisch. \\
\texttt{japanese} & L+R & \texttt{\string\JP} & Japanisch; automatisches Routing erhält die Priorität gemeinsamer Han-Zeichen. \\
\texttt{wenjin} & L+R & \texttt{\string\WJ} & WenJin Mincho mit CJK-Fallback-Kette. \\
\texttt{shanggu} & G & --- & Global orientierte CJK-Familie für traditionelles Chinesisch. \\
\texttt{sim} & G & --- & System-CJK-Familie für vereinfachtes Chinesisch. \\
\texttt{korean} & L+R & \texttt{\string\KR} & Koreanische/Hangul-Familie. \\
\texttt{tangut} & L+R & \texttt{\string\TG} & Tangut mit CJK-bewusstem automatischem Routing. \\

\texttt{mongolian} & L & \texttt{\string\MO} & Mongolisch mit eingebauter vertikaler Layout-Route. \\
\texttt{mongolian\_baiti} & L & \texttt{\string\MOb} & Mongolian Baiti mit vertikaler Layout-Route. \\
\texttt{manchu} & L & \texttt{\string\MC} & Mandschurisch mit vertikaler Layout-Route. \\
\texttt{turkic} & L & \texttt{\string\OT} & Alttürkisch. \\
\texttt{uyghur} & L & \texttt{\string\UY} & Altuigurisch mit vertikaler/gedrehter Layout-Route. \\

\texttt{vietnamese\_quocngu} & L & \texttt{\string\VI} & Vietnamesisches Quốc Ngữ. \\
\texttt{vietnamese\_hannom} & L & \texttt{\string\HN} & Vietnamesisches Hán-Nôm / CJK-Text. \\
\bottomrule
\end{longtable}

\subsection{Erweiterungsschnittstellen für Fortgeschrittene}

Diese Befehle sind für die Erweiterung von IMPE gedacht und nicht für die
gewöhnliche Dokumentkonfiguration.

\begin{longtable}{@{}p{0.39\linewidth}p{0.55\linewidth}@{}}
\toprule
\textbf{Befehl} & \textbf{Zweck} \\
\midrule
\endfirsthead
\toprule
\textbf{Befehl} & \textbf{Zweck} \\
\midrule
\endhead
\texttt{\string\FontRegisterFamily\{...\}} &
Registriert eine neue Schriftfamilie und ihre lokalen/globalen Metadaten. \\
\texttt{\string\LayoutPresetRegister\{...\}} &
Registriert eine strukturierte Layout-Voreinstellung aus internen Layout-Komponenten. \\
\texttt{\string\LayoutPresetDeclare\{...\}} &
Niedrigere Schnittstelle zur Deklaration von Layout-Voreinstellungen, die beim Aufbau des Systems verwendet wird. \\
\bottomrule
\end{longtable}

Die vollständige Registrierungssyntax und die Erweiterungsregeln sind in
\texttt{docs/FONTS.md}, \texttt{docs/LAYOUTS.md} und
\texttt{docs/SYSTEM.md} beschrieben.

\section{Showcase}\label{app:showcase}

Die folgenden Seiten geben den vollständigen kanonischen IMPE-Showcase wieder.
Der Showcase ist kein zweites Referenzhandbuch, sondern ein visueller Begleiter,
der mehrsprachiges Font-Routing, komplexes Shaping, rechts-nach-links laufende
Schriften, vertikales Schreiben, regionale CJK-Formen und Kombinationen von
Dokumentfeatures demonstriert.

Der neueste Showcase kann über den
\href{\IMPEShowcaseURL}{aktuellen main-Branch} bezogen werden; die zu diesem
Handbuch gehörige Ausgabe wird durch den
\href{\IMPETaggedShowcaseURL}{Tag v\IMPEVersion{}} bewahrt.

\IfFileExists{\IMPEShowcaseFile}
  {\includepdf[pages=1-3,pagecommand={}]{\IMPEShowcaseFile}%
   \includepdf[pages=4-6,landscape=true,pagecommand={}]{\IMPEShowcaseFile}%
   \includepdf[pages=7-,pagecommand={}]{\IMPEShowcaseFile}}
  {\PackageError{impe-manual}
     {Erforderliche Repository-Ressource ../../_showcase/main.pdf fehlt}
     {Aus doc/de kompilieren oder scripts/build_manual.ps1 aus dem Repository-Stammverzeichnis ausfuehren.}}

\end{document}
